\section{术语消化：本期关键词索引}

\begin{longtable}{p{0.24\textwidth}p{0.34\textwidth}p{0.33\textwidth}}
\toprule
术语 & 一句话解释 & 在本期中的作用 \\
\midrule
Robotics & 机器人学，研究感知、控制、规划和实体执行 & 全片主题。 \\
Embodied AI & 具身智能，智能体在环境中感知和行动 & 中文语境常用，但谭杰更倾向说机器人。 \\
Sim2Real & 从仿真训练迁移到真实世界 & 谭杰早期范式转移的核心。 \\
Reinforcement Learning & 强化学习，通过 reward 试错优化策略 & 解决 locomotion 和控制的重要方法。 \\
PPO & Proximal Policy Optimization，常用强化学习算法 & 谭杰早期工作采用的技术背景之一。 \\
MPC & Model Predictive Control，模型预测控制 & 传统机器人控制路线代表。 \\
VLM & Vision-Language Model，视觉语言模型 & 能看图和理解语言，但通常不直接输出动作。 \\
VLA & Vision-Language-Action，视觉语言动作模型 & 机器人基座模型的重要形式。 \\
Motion Transfer & 动作迁移，把动作能力迁到不同本体 & Gemini Robotics 1.5 的关键机制之一。 \\
World Model & 预测世界状态变化的模型 & 帮助机器人规划动作后果。 \\
V-L-V & Vision-Language-Vision & 谭杰描述世界模型的一种方式。 \\
Synthetic Data & 合成数据，通常来自仿真或程序生成 & 弥补真实机器人数据不足。 \\
Tactile Sensing & 触觉感知，测量接触力、滑动和材质 & 灵巧手和精细操作的重要输入。 \\
Dexterous Hand & 灵巧手，多指复杂操作本体 & 触觉和 manipulation 的难点集中区。 \\
Generalist & 通用型机器人，可跨任务泛化 & 长期目标，会压缩 specialist。 \\
Specialist & 专用型机器人，解决特定垂直任务 & 短期更容易落地。 \\
\bottomrule
\end{longtable}

\subsection{本章小结}

本期术语表明，机器人不是单个模型名，而是一组横跨仿真、动作、世界模型、触觉、硬件和安全的系统问题。理解术语之间的关系，比追逐单个 demo 更重要。

